%% ma: L11-B15-0074
%% lop: 11
%% chuong: 5
%% bai: 15
%% dang: 11.15.5
%% loai: DS
%% muc: [NB, TH, TH, VDC]
%% dap_an: "ĐSĐS"
%% nguon: "tai-lieu-hoc-tap-mon-toan-lop-11-theo-chuong-trinh-giao-duc-pho-thong-2018.pdf (Lớp Toán Thầy Duy, Phạm Lê Duy, THPT Chu Văn An), Chương 3 Bài 1, Phần 2, Câu 4 (tr.337)"
%% kien_thuc: "Lớp 11 Bài 15 (giới hạn của dãy số)"
%% trang_thai: cho-duyet
%% ly_do: "Dạng mới đề xuất (11.15.x), chờ giáo viên duyệt dạng; tài liệu gốc không có lời giải mẫu, đáp án do kho tự giải và kiểm bằng sympy."
%% ghi_chu: "Ý d) của đề gốc không nêu rõ cách xác định b, c; kho hiểu b√3 + c là cách viết với c = 0 (đáp án gốc: Sai). Nhờ giáo viên xem lại cách phát biểu."
\begin{ex}
Cho giới hạn $L=\lim\dfrac{\sqrt[3]{an^3+5n^2-7}}{\sqrt{3n^2-n+2}}$. Khi đó
\choiceTF
{\True Khi $a=1$ thì $L=\dfrac1{\sqrt3}$.}
{Khi $a=0$ thì $L=\dfrac13$.}
{\True Khi $a>0$ thì $L>0$.}
{Khi $L=b\sqrt3+c$ với $b,c$ là các tham số thì $P=\dfrac{a+c}{b^3}=\dfrac12$.}
\loigiai{\sloppy Với $a\ne0$: \[L=\lim\dfrac{\sqrt[3]{a+\frac5n-\frac7{n^3}}}{\sqrt{3-\frac1n+\frac2{n^2}}}=\dfrac{\sqrt[3]a}{\sqrt3}.\]
a) $a=1$: $L=\dfrac1{\sqrt3}$. Đúng.
b) $a=0$: tử $=\sqrt[3]{5n^2-7}\sim n^{2/3}$, mẫu $\sim\sqrt3\,n$, nên $L=0\ne\dfrac13$. Sai.
c) $a>0$: $L=\dfrac{\sqrt[3]a}{\sqrt3}>0$. Đúng.
d) $L=\dfrac{\sqrt[3]a}{3}\sqrt3$ nên $b=\dfrac{\sqrt[3]a}3$, $c=0$, và $P=\dfrac{a}{b^3}=27\ne\dfrac12$. Sai.}
\end{ex}
